\pagestyle{highlightmain}
\highlightbar{
    \section{SCFI}
    \textit{09/2012 - 09/2013}\\
    \\SCFI est une PME française de services informatiques spécialisée dans le développement web et les solutions numériques pour les entreprises.
    \\Elle propose des prestations de développement applicatif, d'intégration de systèmes et de maintenance aux petites et moyennes structures.
    \vspace{0.5em}\\
    \location{Paris, France}
    \homepage{scfi.fr}{https://www.scfi.fr/}
    \linkedin{SCFI}{https://www.linkedin.com/company/scfi}
    \vspace{0.5em}\\
    Référent :\\
    \contact{Christophe TATIN, CTO}{christophe.tatin@scfi.fr}
}
\mainbar{
    \section[\faGears]{SCFI - Développeur et intégrateur web}
    \textit{09/2012 - 09/2013}

    Contributions techniques à divers projets web en alternance : sites institutionnels, applications métier internes et plateformes e-commerce.

    \subsection{Principales responsabilités}
    \begin{itemize}
        \item Contribution à la maintenance et à l'évolution de l'intranet d'une association d'aide à domicile : planification, facturation, paie et processus opérationnels.
        \\\item Participation aux réunions régulières avec le client pour recueillir ses retours, analyser ses besoins opérationnels et adapter les fonctionnalités.
        \\\item Maintenance et amélioration du CMS interne de l'entreprise et de ses outils d'administration.
        \\\item Contribution à la maintenance et aux évolutions fonctionnelles de plateformes e-commerce et de sites institutionnels.
        \\\item Développement de composants web front-end et back-end avec \textcolor{accent}{ASP}, \textcolor{accent}{PHP}, JavaScript, \textcolor{accent}{jQuery}, AJAX, HTML et CSS.
        \\\item Intervention sur le déploiement et l'hébergement dans des environnements \textcolor{accent}{IIS} et Linux.
    \end{itemize}

    \subsection{Technologies \& outils}
    \tag{ASP} \tag{PHP} \tag{JavaScript} \tag{jQuery} \tag{AJAX}\\
    \tag{HTML} \tag{CSS} \tag{jQuery UI}\\
    \tag{Dreamweaver} \tag{Photoshop}\\
    \tag{IIS} \tag{Linux}

    \subsection{Langues utilisées}
    \tag{Français}[quotidien, langue maternelle]
}
\makebody
\clearpage